% =====================================================================
% §2.6 — Dal Grain all'audio (sola prosa + diagramma d'architettura)
%
% Pattern: risposta al ponte di §2.5 -> Grain come contratto ->
%          renderer disaccoppiati -> STEMS/MIX -> cache SHA-256 ->
%          export REAPER -> tesi del loop abitabile -> LSP ->
%          diagramma (alla Truax: dopo i mattoni) -> ponte a §3
%
% TODO prima della submission:
%   [ ] footnote versioni software (richiesta interna, mail 2026-05):
%       python --version; pip show numpy soundfile; csound --version
%   [ ] produrre il diagramma (figures/architettura.pdf) dalla pipeline:
%       YAML -> parser/IR -> loop -> Grain -> {NumPy | Csound} ->
%       {STEMS | MIX} -> .aif (+ .rpp, + partitura) — piccolo, una colonna
%   [ ] verificare coerenza terminologica: "stem" vs "stream" nel testo
% =====================================================================

% ----------------------------------------------------------------
\subsection{Dal \texttt{Grain} all'audio}\label{sec:render}

Il percorso inverso parte dal punto in cui il cappello di
\S\ref{sec:architettura} si era fermato: il loop di generazione campiona
la IR e materializza la lista di \texttt{Grain}. Quella lista è il
contratto fra generazione e uscita: tutto ciò che segue la consuma senza
poterla modificare.

A consumarla è un renderer intercambiabile dietro un'interfaccia comune:
quello nativo esegue l'overlap-add direttamente in NumPy, senza
dipendenze esterne; un adapter Csound è disponibile per chi lavora in
quella tradizione, e l'aggiunta di altri back-end non richiede modifiche
al resto del sistema.\footnote{Ambiente di riferimento: Python [TODO],
NumPy [TODO], Csound [TODO].} Ortogonale alla scelta del renderer è la
modalità di uscita: un file unico (\emph{mix}) oppure un file per stream
(\emph{stems}), con gli onset relativi al proprio stream.

È la seconda modalità ad abilitare il workflow su cui il sistema
scommette. Ogni stream riceve un fingerprint SHA-256 calcolato sul
proprio dizionario YAML: alla build successiva vengono ri-renderizzati i
soli stream il cui fingerprint è cambiato, e gli stem orfani di stream
rimossi o rinominati sono eliminati automaticamente. Accanto all'audio
il sistema può esportare un progetto REAPER: una traccia per stream, con
l'item posizionato sulla timeline secondo l'onset e la durata dichiarati
--- la struttura temporale dello YAML trasferita in DAW senza
ricostruzione manuale. Il costo di una modifica torna così proporzionale
alla modifica: su un brano di decine di stream, toccare uno stream
costa uno stream. È questa catena --- fingerprint, rendering selettivo,
progetto che si riapre già montato --- a tenere il ciclo
scrivi--renderizza--ascolta abbastanza corto da poterlo abitare; e
all'altro capo del ciclo, prima ancora del rendering, un Language Server
valida la specifica durante la scrittura.

\begin{figure}[t]
  \centering
  \begin{tikzpicture}[
      node distance=4pt and 4pt,
      every node/.style={font=\fontsize{7}{8}\selectfont},
      box/.style={draw, rectangle, minimum height=4mm, inner sep=2pt,
                  align=center, font=\fontsize{7}{8}\ttfamily},
      lbl/.style={align=center, font=\fontsize{7}{8}\ttfamily},
      arr/.style={-{Latex[length=1.4mm]}, thick}
    ]
\node[box] (yaml) {YAML};
\node[box, right=10pt of yaml] (gen) {Generator};
\node[box, below=10pt of gen] (ir)
      {Stream / VoiceManager / Controller$\times$4};
\node[lbl, below=5pt of ir] (irlbl)
      {List[List[Grain]]};
\node[box, below=5pt of irlbl] (renderer)
      {AudioRenderer (ABC)};
\node[box, below=5pt of renderer] (np)
      {Numpy\\AudioRenderer};
\node[lbl, below=5pt of np] (aif1) {.aif};
\node[box, right=5pt of renderer, align=center] (cache)
      {Stream\\CacheManager\\\fontsize{6}{7}\selectfont SHA-256\\\fontsize{6}{7}\selectfont (STEMS)};
\node[box, left= 5pt of renderer, align=center] (rp)
      {ReaperProject\\Writer};
\node[lbl, below=10pt of rp] (rpp) {.rpp};
\draw[arr] (yaml) -- (gen);
\draw[arr] (gen)  -- (ir);
\draw[arr] (ir)   -- (irlbl);
\draw[arr] (irlbl)-- (renderer);
\draw[arr] (renderer) -- (np);
\draw[arr] (np) -- (aif1);
%\draw[arr] (ir) to[bend right=15] (rp.north);
\draw[arr, dashed] (ir) to (rp.north);
\draw[arr, dashed] (rp) -- (rpp);
\draw[arr, dashed] (aif1.west) to[bend left=25] (rp.south east);
\draw[arr, dashed] (cache) to[bend left=15] (np);
\end{tikzpicture}

\caption{La pipeline: dalla specifica YAML alla IR, dal
    loop di generazione ai \texttt{Grain}, dai renderer intercambiabili
    agli stem (con cache per-stream), al progetto DAW e alla partitura.}\label{fig:arch}
\end{figure}

C'è un ultimo prodotto di ogni rendering, finora usato senza essere
tematizzato: accanto all'audio, il sistema genera la partitura grafica
che ha accompagnato gli esempi di questa sezione. È l'altro capo del
ciclo --- ciò che si guarda fra un ascolto e l'altro --- ed è l'oggetto
della sezione che segue.


